\section{Verification purpose}

This case verifies that the standard ELMFIRE executable initializes, schedules, transports, and merges multiple point ignitions using only prescribed physical inputs and observable model outputs. It preserves the scientific intent of Adam Laird's dual-ignition, staggered spot-merging, and near-boundary component scenarios.

\section{Mathematical formulation and numerical implementation}

Each ignition introduces a local negative level-set region at position $(x_i,y_i)$ and time $t_i$. ELMFIRE then advances $\phi$ according to $\partial_t\phi+V_n|\nabla\phi|=0$. With homogeneous fuel, zero slope, and zero wind, simultaneous mirrored sources must remain reflection symmetric. Adding uniform wind must produce a measurably different arrival-time solution. Delayed sources must appear at their requested times, and clipping at a domain boundary must not suppress or wrap a source.

\section{Derivation of expected behavior}

The no-wind forcing and ignition pair are invariant under reflection about their mapped raster axes; therefore the discrete arrival-time solution has the same symmetries. Because ELMFIRE maps coordinates to one-based cells with a ceiling operation, the evaluator reflects the largest balanced row window about the actual ignition row rather than an intercell domain line. The wind simulation condition changes only the wind rasters, so a nonzero solution difference is a metamorphic check of weather coupling. For scheduled ignition $i$, the containing raster cell should receive an arrival time within one \OutputInterval{} of $t_i$. The eight boundary sources are placed one five-cell margin inside the square domain and all must initialize.

% BEGIN EXPLICIT METRIC DERIVATION
For any two arrival-time arrays \(A\) and \(B\), the implemented comparison is
\[
 E(A,B)=\frac{\operatorname{mean}_{i\in V}|A_i-B_i|}
 {\max[P_{90}(A_i:i\in V),1~\mathrm{s}]},
 \qquad V=\{i:A_i,B_i\text{ are finite}\}.
\]
For the two simultaneous ignitions at \((-52.5,0)\) and \((52.5,0)\) m,
the reflection reference is exact symmetry, \(E=0\); the metric is the
larger of the left--right and centered top--bottom reflection errors and
must satisfy \(E\leq0.02\).  The \(5~\mathrm{mph}\), 225-degree wind run
is compared with the no-wind run by the same formula and must give
\(E\geq0.05\), demonstrating a resolved forcing response.
For scheduled ignitions \((x,y,t)\) equal to
\((0,600,0)\), \((0,300,1200)\), \((0,0,2400)\),
\((0,-300,3600)\), and \((0,-600,4800)\), the timing metric uses the
simulated time of arrival (TOA) at each ignition location:
\[
 E_t=\max_j|TOA(x_j,y_j)-t_j|\leq1200~\mathrm{s}.
\]
Finally, the count metric is
\(\sum_{j=1}^{8}\mathbf{1}[TOA(x_j,y_j)\leq1200~\mathrm{s}]=8\) at the
eight generated edge/corner points \((\pm975,0)\), \((0,\pm975)\), and
\((\pm975,\pm975)\) m.
% END EXPLICIT METRIC DERIVATION

\section{Verification metrics and acceptance criteria}

Reflection error is the larger normalized mean absolute difference between the no-wind arrival field and its reflection about each mapped ignition axis; it must not exceed 0.02. Wind response is the normalized mean absolute difference between the wind and no-wind fields and must be at least 0.05. The maximum scheduled-source timing error must not exceed one output interval. All eight boundary-source cells must have finite arrival times no later than one output interval. These metrics were selected because they test invariants of the coupled ignition/level-set implementation without reproducing ELMFIRE's internal algorithms.

\section{Simulation configuration and predicted outcome}

All simulation conditions use a \GridSpacing{} square grid in EPSG:32610, Fuel Model 10, 3/5/6 percent dead-fuel moisture, and flat terrain. The simulation conditions contain two simultaneous no-wind sources, the same pair under 5 mph wind from 225 degrees, five staggered sources, and eight edge/corner sources. In the staggered simulation condition the first source is anchored at (t=0) and the other four follow at 1200-s intervals. This common translation of the original absolute schedule preserves every relative delay while preventing versions that have no active front from advancing past a strictly future first ignition. The predicted fields are respectively symmetric merger, wind-distorted merger, correctly timed merger, and eight clipped but valid boundary fires.

\begin{center}\begin{tabular}{p{0.22\textwidth}p{0.22\textwidth}p{0.12\textwidth}p{0.34\textwidth}}\toprule
Parameter & Configured value & Units & Role \\
\midrule Domain/grid & 2000 m square; 5 & m/cell & Resolves source placement and merging \\
Fuel/moisture & FM10; 3/5/6 & --- / percent & Homogeneous surface response \\
Time control & step 1, maximum 5; stop 14400 & s & Exact timestep multiple for scheduled sources \\
Simulation conditions & dual calm, dual wind, staggered, boundary & --- & Exercises symmetry, weather, timing, and clipping \\
Output cadence & 1200 & s & Timing-error acceptance scale \\
\bottomrule\end{tabular}\end{center}

Figure~\ref{fig:input-configuration} shows the prescribed spatial condition used in this experiment. The displayed field represents the initial level-set geometry or a controlling fuel, terrain, moisture, wind, or canopy quantity, as identified by the figure labels. The plotted values come from the prepared simulation inputs, not from the expected response.

\begin{figure}[H]
  \centering
  \IfFileExists{../figures/input_configuration.pdf}{%
    \includegraphics[width=0.98\textwidth,height=0.46\textheight,keepaspectratio]{input_configuration.pdf}%
  }{\fbox{\parbox{0.88\textwidth}{Prepared-input evidence is unavailable. The prescribed configuration is described in the text; this absence does not establish that it was simulated.}}}
  \caption{Prescribed initial level-set geometry for two simultaneous ignitions without wind. Coordinates are grid-cell indices on the 5 m grid; the field is dimensionless.}
  \label{fig:input-configuration}
\end{figure}

\section{Actual simulation results}

The generated result is \textbf{\VerificationStatus}; \PassingMetricCount{} of \MetricCount{} required metrics pass across \VariantCount{} simulation conditions. Missing ELMFIRE output yields NOT EVALUABLE rather than a pass.

\begin{figure}[H]\centering
\IfFileExists{../figures/verification_summary.pdf}{\includegraphics[width=\textwidth]{verification_summary.pdf}}{\fbox{No case-specific result figure is available.}}
\caption{Case-specific time-of-arrival fields for the four ignition-topology simulation conditions.}
\end{figure}

\section{Calculated metrics and verification decision}

\begin{center}\begin{tabular}{@{}llll@{}}\toprule
Metric & Expected & Calculated & Status \\\midrule
\MetricRows
\bottomrule\end{tabular}\end{center}
The authoritative decision is \textbf{\VerificationStatus} from calculated verification and validation metrics.

\section{Technical interpretation}

A symmetry failure points to directional indexing or level-set transport. No wind response indicates that the weather field did not reach fire behavior. Timing error implicates the scheduled-ignition path. Missing edge ignitions indicates boundary-index clipping or coordinate conversion failure.

\section{Scope and limitations}

This is deterministic implementation verification, not fire-behavior validation. It does not expose the internal Fortran ignition calculations directly; it checks their observable behavior through an unmodified ELMFIRE executable.
